\section{Overview}\label{1.1}

Large language models have made it easy to get a fluent answer to almost any
question. What they have not made easy is knowing whether that answer is
\emph{correct}. A model with a few billion parameters, small enough to run
on a laptop GPU with no internet connection, will answer a biology question
just as confidently whether it is right or wrong. For a student using the
system to revise for an exam, that is a serious problem: a wrong answer
delivered with total confidence is worse than no answer at all, because
there is nothing to signal that it needs checking.

\textbf{AURA-Bio} (Adaptive Utility-Routed Assistant for Biology) is a
response to that problem. It is a local-first tutoring system for biology
that treats a small open-source language model not as the final word, but as
one component inside a larger pipeline: the question is understood and
routed, relevant evidence is retrieved from a large, verified biology
corpus, an answer is drafted against that evidence, and the answer is then
independently checked for factual support, correct terminology, sound
process ordering, and correct arithmetic before it is shown to the
learner. Every one of those stages is logged, so any answer can be traced
back to exactly what evidence it used and what checks it passed or failed.

Biology was chosen as the first subject not because the architecture is
biology-specific, but because it is broad enough to exercise every part of
the system: factual recall, multi-step processes, quantitative genetics,
terminology, and material a small model can have large information about
already as well as leaving gaps for the pipeline to improve the outputs upon while
remaining a single, well-scoped domain for a first implementation. The same
architecture is designed from the outset to extend to other technical
subjects. Chapter~5 sets out concretely how.

\subsection{What Kind of System This Is}\label{1.1.1}

It is worth being precise about what AURA-Bio is not. It is not a general
chatbot with a biology theme, and it is not a wrapper that simply pastes
retrieved text in front of a language model and hopes for the best (a
pattern usually called Retrieval-Augmented Generation, or RAG). The
difference is the second half of the pipeline. A plain RAG system stops
once it has generated an answer from retrieved context. AURA-Bio continues,
the generated answer is verified against the evidence it was given,
checked by tools that recompute arithmetic and re-derive genetic ratios
rather than merely reading the model's stated answer, and, atleast during
development, reviewed by an independent critic model before the design is
trusted. A conventional RAG pipeline looks like this:

\begin{center}
\ttfamily
Question $\rightarrow$ Retrieve documents $\rightarrow$ Generate answer
\end{center}

The journey of a single question through AURA-Bio, for a reader who wants
the whole picture at a glance before the detailed treatment in
Chapter~\ref{3}, looks like this:

% Wrapped in a minipage so the diagram is never split across a page break:
% a minipage is an unbreakable box, so LaTeX moves the whole thing to the
% next page rather than cutting it in half.
\begin{center}
\begin{minipage}{\textwidth}
\centering
\footnotesize
\begin{verbatim}
+----------------------------------------------------------+
|                      USER'S QUESTION                     |
+----------------------------------------------------------+
                             |
                             v
+----------------------------------------------------------+
|  ROUTE                                                   |
|    task type . subdomain . learner level . teaching mode |
+----------------------------------------------------------+
                             |
                             v
+----------------------------------------------------------+
|  RETRIEVE  (in parallel)                                 |
|    sparse : BM25 keyword search over the corpus          |
|    dense  : BGE semantic vector search                   |
+----------------------------------------------------------+
                             |
                             v
+----------------------------------------------------------+
|  FUSE (RRF)  ->  RERANK (cross-encoder)                  |
+----------------------------------------------------------+
                             |
                             v
+----------------------------------------------------------+
|  SELECT EVIDENCE   (per-source cap, prompt budget)       |
+----------------------------------------------------------+
                             |
                             v
+----------------------------------------------------------+
|  DRAFT ANSWER   (local 4B model, on-device)              |
+----------------------------------------------------------+
                             |
                             v
+----------------------------------------------------------+
|  VERIFY                                                  |
|    claim support . terminology . process order .         |
|    arithmetic / units / pathways . entity authority      |
+----------------------------------------------------------+
                |
                v
        defect found? --- yes --> REVISE, then re-VERIFY
                |                      (bounded to 2 passes)
                no
                |
                v
+----------------------------------------------------------+
|  FINAL ANSWER  --  citations, confidence, full trace     |
+----------------------------------------------------------+
\end{verbatim}
\end{minipage}
\end{center}

\subsection{Why This Matters for Small Models Specifically}\label{1.1.2}

The largest commercial language models are, at the time of writing, quite
good at biology out of the box. The reason a project like AURA-Bio matters
is that not every learner, school, or institution can or should depend on a
model that only runs on someone else's server, at someone else's cost, on
someone else's terms of data handling. A model small enough to run on a
four-year-old laptop GPU is the only kind of model many settings can
actually deploy. The question this project asks is therefore precise:
\emph{given a small model that will make mistakes, what system built around
it makes those mistakes visible, checkable, and where possible, corrected
without needing a bigger model?}

\section{The Problem Being Addressed}\label{1.2}

Small open-source language models are attractive because they can run
locally, preserve privacy, remove recurring API costs, and function without
permanent dependence on a commercial platform. In practice, however, they
exhibit a recognisable set of failure modes:

\begin{enumerate}
    \item They may hallucinate facts, stating something false with the same
        confidence as something true.
    \item Their biology knowledge may be incomplete, outdated, or simply
        absent for less common topics.
    \item They may confuse similar biological terms like mitosis for meiosis,
        an artery for a vein, transcription for translation.
    \item They may produce explanations pitched at the wrong level: too
        simple for an undergraduate, too technical for a school student.
    \item They may fail to adapt their teaching style to what the learner
        actually asked for.
    \item They may produce answers with no indication of where the
        information came from.
    \item They may misrepresent a multi-step biological process, reversing
        or omitting steps.
    \item They may fail to communicate uncertainty, presenting a debated or
        incomplete finding as settled fact.
    \item They may be fluent and well-structured while being scientifically
        wrong, which is, if anything, more dangerous than an answer that
        reads as unconfident.
    \item Claims that a system has \emph{improved} a model are often made
        without any measurable evidence to support them.
\end{enumerate}

AURA-Bio addresses these limitations not by attempting to make the
underlying model itself smarter, but by surrounding it with a structured
reasoning and verification pipeline that catches what the model gets wrong
and is transparent about what it does not know.

\section{Core Idea}\label{1.3}

The central idea investigated in this project is that:

\begin{quote}
\emph{A small local language model can become significantly more useful for
biology education when combined with domain-specific knowledge retrieval,
adaptive teaching, evidence checking, structured critique, and measurable
evaluation.}
\end{quote}

Rather than treating the model as a standalone chatbot, AURA-Bio treats it
as one component within a larger system that, for every question:
understands what is being asked, determines the biology subdomain and the
learner's academic level, chooses an appropriate teaching style, retrieves
relevant evidence, reranks and filters that evidence, generates a draft
answer, verifies its claims and terminology, uses a critic to identify
weaknesses during development, revises the answer where a genuine defect is
found, and returns the final answer with citations and a confidence signal, and
logs the entire pipeline for later analysis and, eventually, for training
data.

\section{Objectives}\label{1.4}

\subsection{Primary Objectives}\label{1.4.1}

\begin{enumerate}
    \item Develop a local-first biology tutoring system using an
        open-source language model.
    \item Rebuild a validated biology corpus (a raw-data ceiling of
        approximately 50\,GB was set, though quality took priority over
        reaching it) for retrieval and possible future domain adaptation.
    \item Implement document ingestion, cleaning, deduplication, chunking,
        and metadata generation.
    \item Implement dense semantic retrieval using BGE embeddings and
        sparse keyword retrieval using BM25, combined through hybrid
        fusion.
    \item Implement reranking to select the most useful retrieved context.
    \item Build a biology-specific task, subdomain, learner-level, and
        teaching-style router.
    \item Verify biological claims against retrieved evidence, and detect
        unsupported statements, terminology mistakes, and process
        inconsistencies.
    \item Use Python-based deterministic tools (SymPy, Biopython, and
        others) for claims that can be checked by recomputation rather than
        by pattern matching.
    \item Use an optional external model during development for critique
        and independent judging.
    \item Implement a structured critique-revision loop, bounded so it
        cannot run indefinitely.
    \item Expose the entire pipeline through trace logs.
    \item Evaluate the contribution of each system component using
        ablation studies measuring what each piece actually buys, rather
        than assuming it helps.
    \item Build a failure-learning dataset for future fine-tuning.
    \item Design the architecture so additional subject domains can later
        be added as modules.
\end{enumerate}

\subsection{Secondary Objectives}\label{1.4.2}

\begin{itemize}
\item Generating revision notes, quizzes, and flashcards.
\item Identifying learner misconceptions.
\item Supporting exam-oriented and research style explanations.
\item Adapting explanations across school and undergraduate levels.
\end{itemize}

\section{Scope}\label{1.5}

\subsection{Initial Scope}\label{1.5.1}

The first implementation focuses on biology, spanning general biology, cell
biology, molecular biology, genetics, evolution, ecology, botany, zoology,
microbiology, human physiology, anatomy, immunology, biochemistry,
biotechnology, developmental biology, neuroscience basics, bioinformatics
concepts, and introductory research-paper interpretation: eighteen
subdomains in total, each represented in the evaluation benchmark described
in Chapter~3.

\subsection{Out of Scope for This Version}\label{1.5.2}

Diagnosis of disease, treatment recommendation, prescription advice, and
clinical decision support are explicitly out of scope. Medical and
health-related responses are treated as educational content only, and are
never presented as professional medical advice. Full coding assistance,
advanced symbolic mathematics, and autonomous scientific research are also
out of scope for this first version, though Chapter~5 discusses how the
same architecture could later support them as separate modules.

\section{Intended Users}\label{1.6}

The system is designed around four user groups. \textbf{School students}
studying biology at secondary or senior-secondary level, who need concept
explanations, exam preparation, homework support, and misconception
correction. \textbf{Undergraduate students} in biology, biotechnology, or
life sciences, who need deeper mechanistic explanations, pathway detail,
and help understanding research papers. \textbf{Teachers and teaching
assistants}, who can use the system to generate question sets, prepare
revision material, and identify common misconceptions. And
\textbf{researchers or lab students}, who may use it to summarise papers,
retrieve information from a controlled corpus, and navigate unfamiliar
terminology.

\section{Report Organisation}\label{1.7}

Chapter~2 surveys the relevant background: retrieval-augmented
generation, hallucination in language models, and the case for
deterministic verification over using a second model as judge. Chapter~3
describes the methodology: the corpus, the retrieval architecture, the
routing system, the verification framework, the critic, and the
benchmark used to measure everything. Chapter~4 presents the measured
results, defining every scoring metric precisely before reporting what
each component of the pipeline contributes. Chapter~5 concludes with the
project's current status, the fine-tuning work now under way, and a
concrete plan for generalising the architecture beyond biology.
